Protein language models (pLMs) have revolutionized our ability to predict protein structures and functions, yet the internal representations they learn remain largely opaque. While existing mechanistic interpretability studies often focus on mapping model features to known biological annotations, the true potential of these models lies in their ability to capture biological principles that have yet to be formally curated. In this work, we present a pipeline that treats pLMs as active generators of scientific hypotheses. By applying sparse autoencoders (SAEs) to the latent representations of \esm, we identify monosemantic features with highly specific activation patterns that do not align with existing \swissprot annotations. We then leverage large language models (LLMs) to semantically interpret the sequence contexts of these "orphan" features, generating explicit, biologically plausible hypotheses for uncharacterized motifs. Our results identify several novel candidates, such as "Charged Turn Motifs" and "Positively-Charged Helix Caps," which exhibit consistent biophysical properties across diverse protein families. This work demonstrates the viability of using foundation models as engines for biological discovery, moving beyond prediction toward the mechanistic understanding of novel biological phenomena.
